\documentclass[aspectratio=169, 10pt]{beamer}
\usetheme{metropolis}

% --- Edinburgh colours ----------------------------------------
\definecolor{UoERed}{HTML}{7A2318}
\definecolor{UoEGold}{HTML}{B8860B}
\definecolor{CodeBG}{HTML}{F7F3ED}
\definecolor{CodeFrame}{HTML}{D5C9B1}
\definecolor{UoEBlue}{HTML}{2a78d6}

\setbeamercolor{palette primary}{bg=UoERed, fg=white}
\setbeamercolor{frametitle}{bg=UoERed, fg=white}
\setbeamercolor{title separator}{fg=UoEGold}
\setbeamercolor{progress bar}{fg=UoEGold, bg=UoEGold!20}
\setbeamercolor{alerted text}{fg=UoERed}
\setbeamercolor{block title}{bg=UoERed!15, fg=UoERed}
\setbeamercolor{block body}{bg=UoERed!5}

% --- Packages -------------------------------------------------
\usepackage[T1]{fontenc}
\usepackage{lmodern}
\usepackage{amsmath, amssymb, mathtools}
\usepackage{listings}
\usepackage{graphicx, booktabs}
\usepackage{tikz}
\usetikzlibrary{arrows.meta, positioning, calc, decorations.pathreplacing}
\usepackage{hyperref}
\hypersetup{colorlinks=true, linkcolor=UoERed, urlcolor=UoERed!70!black}
\setbeamertemplate{navigation symbols}{}

\lstset{
  language=Python,
  basicstyle=\ttfamily\scriptsize,
  keywordstyle=\color{UoERed}\bfseries,
  commentstyle=\color{gray}\itshape,
  stringstyle=\color{UoEGold},
  backgroundcolor=\color{CodeBG},
  frame=single,
  rulecolor=\color{CodeFrame},
  numbers=none, breaklines=true, showstringspaces=false, tabsize=4,
  xleftmargin=4pt, xrightmargin=4pt, aboveskip=6pt, belowskip=4pt,
}

\AtBeginSection[]{
  \begin{frame}
    \vfill\centering
    \usebeamerfont{title}\insertsectionhead\par
    \vfill
  \end{frame}
}

\title{Week 6 --- Multiple Time Series \& VAR Models}
\subtitle{Vector Autoregressions, IRFs, and Granger Causality}
\author{Dr Juan Zurita}
\institute{Economic \& Business Forecasting\\
The University of Edinburgh~$\cdot$~School of Economics}
\date{}

\begin{document}
\maketitle


\section{The VAR Model}

\begin{frame}{From Univariate to Multivariate}
Why go multivariate?
\begin{itemize}
    \item Variables \alert{interact}: GDP affects unemployment affects consumption
    \item Omitting relevant variables biases forecasts
    \item We want to trace \textbf{dynamic interactions} (impulse responses)
\end{itemize}
\end{frame}

\begin{frame}{VAR($p$) Specification}
For $\mathbf{y}_t = (y_{1t}, \ldots, y_{nt})'$:
\[
\mathbf{y}_t = \mathbf{c} + \mathbf{A}_1 \mathbf{y}_{t-1} + \cdots + \mathbf{A}_p \mathbf{y}_{t-p} + \mathbf{u}_t
\]
where $\mathbf{u}_t \sim N(\mathbf{0}, \Sigma)$.

\vspace{0.5em}
\begin{itemize}
    \item Each equation is a \textbf{multivariate regression}
    \item Estimated by \textbf{OLS equation by equation}
    \item Parameters: $n^2 p + n$ --- grows fast!
\end{itemize}
\end{frame}

\section{Impulse Response Functions}

\begin{frame}{IRFs}
An IRF traces the effect of a \alert{one-unit shock} in variable $j$ on variable $i$ over time.

\vspace{0.5em}
Write the VAR in MA($\infty$) form:
\[
\mathbf{y}_t = \boldsymbol{\mu} + \sum_{s=0}^{\infty} \Phi_s \mathbf{u}_{t-s}
\]

\textbf{IRF:} $\frac{\partial y_{i,t+h}}{\partial u_{j,t}} = [\Phi_h]_{ij}$

\vspace{0.5em}
\alert{Identification problem:} shocks may be correlated ($\Sigma$ not diagonal). \\
Common solution: \textbf{Cholesky decomposition} of $\Sigma$.
\end{frame}

\section{Granger Causality}

\begin{frame}{Granger Causality}
$x$ \alert{Granger-causes} $y$ if past values of $x$ help predict $y$ beyond $y$'s own past.

\vspace{0.5em}
\textbf{Test:} in the equation for $y_t$, test $H_0$: all coefficients on lagged $x$ are zero.

\vspace{0.5em}
\begin{block}{Caution}
Granger causality $\neq$ true causality. It is about \emph{predictive content}, not causal mechanisms.
\end{block}
\end{frame}

\begin{frame}{VAR Forecasting}
\begin{itemize}
    \item Iterate the VAR forward: $\hat{\mathbf{y}}_{T+h|T} = \mathbf{c} + \sum_{i=1}^{p} \mathbf{A}_i \hat{\mathbf{y}}_{T+h-i|T}$
    \item Lag order selection: AIC, BIC, HQ
    \item Forecast uncertainty: bootstrap or asymptotic intervals
\end{itemize}

\vspace{0.5em}
\textbf{Next week:} state space models and the Kalman filter.
\end{frame}


\end{document}
